\section{执行判断与证据分级}
\subsection{建议保留什么，先停止什么}
Research W 的独立价值不在于把“成像、更新背景、再成像”重新组合。DBIM 已在不均匀当前背景上反复线性化；layer stripping 与 invariant embedding 已研究边界响应的向内推进；Marchenko/redatuming 已研究从单侧反射恢复内部 Green 响应；Bayesian 与 goal-oriented inversion 已把参数不确定性传递到任务量。\cite{L_DBIM,L_Somersalo,L_EM2017,L_Marchenko22,L_Goal}

更值得立项的问题是：
\begin{quote}
在明确的单侧采集、合法材料类、目标任务类和噪声水平下，哪些观测无法区分的环境，会产生不同的完整目标响应？能否识别并保留这些传播分支，同时证明哪些分支可以安全合并？
\end{quote}

这是把关注点从“墙体材料是否重建正确”转移到“对指定成像任务，环境还存在哪些不可消除的传播歧义”。两者可能分离：材料不唯一而任务可辨识，也可能墙体回波拟合很好而目标传播不可靠。

\textbf{首选路线：传播可辨识性与不可能性。} 首轮以普通非磁、被动材料的有限非线性见证、联合目标补偿和任务误差下界为主。概率传播压缩作为其建设性后续；A17 是可选计算接口。不要先训练 Flow/Diffusion，也不要先搭建完整“逐层智能成像系统”。

\subsection{本轮实质结果}
本报告给出如下可检查成果，而非宣称已解决任意复杂墙体成像。
\begin{enumerate}
\item 明确总材料场与“墙/目标”语义分解的区别，给出无结构约束时的标签重分配自由度。
\item 给出任务可辨识的精确纤维判据、局部条件、有限维稳定常数及无限维反例；把目标和仪器补偿纳入局部观测算子。
\item 构造非磁、无损介质墙的精确单频 Maxwell 反例：两墙反射相同，双程目标响应不同；允许目标位移后，完整联合数据又完全相同。
\item 提出应保留的四块传播签名，覆盖背景响应、入射、接收及目标域内部反馈；推导远离共振时的响应误差界。
\item 证明在原联合先验不变的条件下，传播近似如何通过 likelihood 误差控制目标后验误差；指出仅聚类墙体后验不足以保留目标条件分布。
\item 用随机 Schur 响应定义可回退的计算前沿，并说明它通常不是按深度嵌套的物理层。
\item 核验 A17 最新 closure 与效率 follow-up，修正“只有冻结 GN、没有非线性后果”的过时总结；同时保留 A26、A30、F3、P2 的负结果。
\item 实际运行一维精确传输矩阵核验和带宽扫描，附完整代码与原始结果。二维复杂墙体的关键实验尚待执行。
\end{enumerate}

\Needspace{14\baselineskip}
\subsection{数学成立与论文原创性必须分开}
下文采用五类状态。
\begin{longtable}{p{27mm}p{121mm}}
\toprule
状态 & 含义\\
\midrule
已知工具 & resolvent、Schur、Woodbury、T-matrix、条件概率、二点检验与标准伴随残差恒等式。\\
本次可证推论 & 在正文明确假设下给出证明，但尚未建立文献首创性。四块签名稳定界、联合测度不变的压缩界属于此类。\\
解析反例 & 由完整 Maxwell 的严格子问题构造，已数值核对符号与参考平面；适用范围不能外推。\\
待检验机制 & 风险调度、传播分支搜索、可回退前沿及 A17 条件选空间的实际收益。\\
未解决 & 一般未知拓扑墙体的稳定全域后验、可靠跨拓扑混合、同成本算法优势和实际建筑穿透保证。\\
\bottomrule
\end{longtable}

一般的“反射不决定透射”在多模逆散射中已有明确讨论；本报告不会把这一概念、简单共振板反例或通用概率不等式冒称原创 Maxwell 定理。\cite{L_Waagaard} 强 IEEE TAP 贡献仍需把它收紧到可实现材料、有限采集和可计算风险，并用非规则有限尺寸墙体验证。

\subsection{阅读范围与缺口}
本轮核验五个指定仓库的核心决策与报告，固定 commit；A17-X 与 Research F 原 Pro 报告取得原文并检查关键理论节。独立的 L/Z/X 全部 ZIP、P3-L 原报告，以及 A12/A24--A31 每一个独立旧稿并未全部取得。相应讨论明确区分本轮原文证据、实施报告的历史引用和用户任务书，不声称完成未发生的全文阅读。后文的既有研究证据映射详列来源状态。

